\subsection{Scaling Laws and Empirical Observations}
\label{sec:scaling-laws}

A defining feature of modern deep learning---and large language models in
particular---is the remarkable predictability with which performance improves as
a function of scale.  This section surveys the empirical scaling laws that
govern pretraining, post-training via reinforcement learning, and inference-time
computation, and discusses their implications for the GRPO scaling study
undertaken in this capstone.

%% ---------- Classical Scaling Laws ----------
\subsubsection{Classical Scaling Laws for Pretraining}

Kaplan et al.\ introduced the first comprehensive empirical characterisation of
how language model cross-entropy loss depends on three scale
factors~\cite{kaplan2020scaling}.  Training Transformer language models across
more than seven orders of magnitude in parameter count, dataset size, and
compute budget, they established that test loss obeys smooth power-law
relationships with each factor when the other two are not bottlenecked:
\begin{equation}
  L(N) = \left(\frac{N_c}{N}\right)^{\!\alpha_N}\!,\qquad
  L(D) = \left(\frac{D_c}{D}\right)^{\!\alpha_D}\!,\qquad
  L(C) = \left(\frac{C_c}{C}\right)^{\!\alpha_C}\!,
  \label{eq:kaplan-power-laws}
\end{equation}
where $N$ denotes non-embedding parameters, $D$ the number of training tokens,
$C$ the compute budget (in PF-days), and $N_c, D_c, C_c$ are scale constants.
The fitted exponents---$\alpha_N \approx 0.076$, $\alpha_D \approx 0.095$, and
$\alpha_C \approx 0.050$---imply that scaling model size yields larger marginal
returns than scaling data alone.  A key practical corollary was that
compute-efficient training should allocate most of a fixed compute budget to
\emph{larger models trained on relatively modest data}, stopping well before
convergence.  Additional findings included the \emph{universality of
overfitting}: the performance penalty from insufficient data scales predictably
as $N^{0.74}/D$, meaning that an eightfold increase in model size requires only
a roughly fivefold increase in data to maintain the same loss.  Training curves
themselves follow power laws whose parameters are approximately independent of
model size, enabling extrapolation of final loss from early training dynamics.
These scaling laws provided the intellectual foundation for the GPT-3 scaling
strategy and fundamentally shaped how the field reasons about resource
allocation in large-scale pretraining.

%% ---------- Chinchilla Scaling ----------
\subsubsection{Compute-Optimal Training: Chinchilla Scaling}

Hoffmann et al.\ revisited the question of optimal resource allocation by
training over 400 language models ranging from 70 million to 16 billion
parameters on 5 to 500 billion tokens~\cite{hoffmann2022chinchilla}.  Their
central finding overturned the Kaplan prescription: for compute-optimal
training, model size and the number of training tokens should be scaled
\emph{equally}---for every doubling of model size, the number of training tokens
should also be doubled.  Formally, they modelled the loss as a joint function of
parameters and data:
\begin{equation}
  L(N, D) = \frac{A}{N^{\alpha}} + \frac{B}{D^{\beta}} + E,
  \label{eq:chinchilla}
\end{equation}
where $E$ represents an irreducible loss term and $\alpha \approx 0.34$,
$\beta \approx 0.28$.  Under a fixed compute budget $C \propto 6ND$ (the
standard FLOPs approximation for Transformer training), the optimal allocation
scales both $N$ and $D$ proportionally with $C$.  This analysis revealed that
many contemporary models---including the original GPT-3 (175B parameters
trained on 300B tokens)---were significantly \emph{undertrained} relative to
their parameter count.  The authors validated their prediction by training
Chinchilla, a 70B-parameter model on $4\times$ more data than the 280B Gopher,
using the same compute budget.  Chinchilla uniformly outperformed Gopher,
GPT-3, and Megatron-Turing NLG (530B) on a wide range of downstream tasks,
demonstrating that compute-optimal data scaling is at least as important as
parameter scaling.

The Chinchilla result has had far-reaching consequences: it established that
the ratio of tokens to parameters is a first-order design variable, shifted
industry practice toward training smaller models on larger corpora (e.g.,
Llama~2 at 70B on 2T tokens), and provided a quantitative framework for
budgeting pretraining runs.  For the present study, the Chinchilla insight is
relevant insofar as the base models used for GRPO fine-tuning---Qwen-2.5 and
Llama-3 families---have been pretrained under approximately compute-optimal
regimes, ensuring that the starting checkpoints represent efficient use of their
respective pretraining budgets.

%% ---------- Scaling Laws for RL Fine-Tuning ----------
\subsubsection{Scaling Laws for Reinforcement Learning Fine-Tuning}

While pretraining scaling laws are well established, the scaling behaviour of
\emph{post-training} via reinforcement learning has only recently begun to
receive systematic study.  RL training differs from pretraining in fundamental
ways: it optimises a reward signal rather than a log-likelihood objective, it
operates on-policy with generated rollouts rather than static data, and it
exhibits saturating rather than power-law performance curves.

Khatri et al.\ conducted the first large-scale systematic investigation of RL
scaling, spanning over 400,000 GPU-hours of
experiments~\cite{khatri2025scaling}.  They proposed a sigmoidal
compute-performance model for RL training:
\begin{equation}
  R_C - R_0 = (A - R_0) \cdot \frac{1}{1 + (C_{\text{mid}} / C)^B},
  \label{eq:rl-sigmoid}
\end{equation}
where $R_C$ is the expected reward at compute budget $C$, $R_0$ is the initial
performance, $A \in [0,1]$ is the asymptotic pass rate (performance ceiling),
$B > 0$ is a scaling exponent governing compute efficiency, and $C_{\text{mid}}$
sets the midpoint of the performance curve.  This sigmoidal form---with its
characteristic slow-start, rapid-improvement, and saturation phases---contrasts
sharply with the unbounded power-law improvement observed in pretraining and
reflects the bounded nature of task-specific accuracy metrics.

Three key findings emerged from this study.  First, \emph{not all RL recipes
reach the same ceiling}: design choices such as loss type and batch size
materially shift the asymptotic performance $A$, whereas interventions commonly
believed to improve peak performance---such as loss aggregation, advantage
normalisation, data curriculum, and length penalties---primarily modulate compute
efficiency $B$ without raising the ceiling.  Second, \emph{methods that appear
superior at small compute budgets can be inferior at scale}, underscoring the
importance of evaluating RL algorithms by their scaling trajectories rather than
fixed-budget snapshots.  Third, \emph{scaling model size substantially improves
both the asymptotic RL performance and downstream generalisation}, confirming
that the quality of the pretrained base model is a dominant factor in RL
outcomes.

Complementary work by Nimmaturi et al.\ derived predictive scaling laws
specifically for GRPO training across Llama~3B/8B and Qwen~3B/7B
models~\cite{nimmaturi2025grposcaling}.  Their analysis confirmed the
exponential saturation pattern and identified that training beyond a certain
fraction of an epoch offers diminishing returns, suggesting that principled
early stopping can reduce GRPO compute costs significantly without sacrificing
downstream performance.

The effect of \emph{group size} $G$---the number of completions sampled per
prompt in GRPO---represents a scaling dimension unique to group-based RL
methods.  Larger groups yield more stable advantage estimates by providing
richer within-group reward statistics, but increase per-step compute linearly.
DeepSeek-R1 used $G = 64$ for its large-scale training, while academic
reproductions have found that values between 8 and 32 offer a practical
trade-off between variance reduction and compute cost.

A particularly consequential finding from DeepSeek-R1 concerns the interaction
between model size and RL training
effectiveness~\cite{deepseek2025r1}.  DeepSeek reported that for smaller models
(1.5B--14B parameters), \emph{distilling} reasoning traces from a larger
RL-trained model into the smaller model via supervised fine-tuning consistently
outperformed applying RL directly to the smaller model.  This suggests that
small models may lack the capacity to independently discover effective reasoning
strategies through RL exploration, even though they can faithfully reproduce
such strategies when provided as training data.  The implication is stark: the
effectiveness of RL fine-tuning is not scale-invariant, and there may exist a
minimum model capacity below which direct RL training yields suboptimal
returns compared to distillation-based approaches.

%% ---------- Test-Time Compute Scaling ----------
\subsubsection{Test-Time Compute Scaling}

A complementary scaling dimension operates not during training but at inference
time.  Snell et al.\ investigated how allocating additional compute at test time
can improve LLM performance on challenging
problems~\cite{snell2024scaling}.  They analysed two primary mechanisms for
scaling test-time computation: (1)~\emph{search against process-based verifier
reward models}, in which multiple candidate solutions are generated and ranked
by a learned verifier; and (2)~\emph{adaptive distribution updating}, in which
the model's response distribution is refined conditioned on the specific prompt
at test time.

Their central finding was that the effectiveness of test-time compute scaling
is \emph{critically prompt-dependent}: on easy problems, additional inference
compute yields minimal gains, while on problems of moderate difficulty---where
the base model achieves non-trivial but imperfect success rates---test-time
scaling can deliver substantial improvements.  This observation motivated a
\emph{compute-optimal} test-time strategy that adaptively allocates inference
compute per prompt based on estimated difficulty, achieving more than $4\times$
improvement in compute efficiency over uniform best-of-$N$ baselines.  Most
strikingly, in a FLOPs-matched comparison, a smaller model equipped with
compute-optimal test-time scaling outperformed a $14\times$ larger model
operating in zero-shot mode, demonstrating that inference-time and
pretraining-time compute are partially substitutable resources.

This result has direct relevance to the long chain-of-thought reasoning
paradigm pioneered by DeepSeek-R1 and OpenAI's o1.  These models effectively
implement a form of test-time compute scaling by generating extended reasoning
traces---sometimes thousands of tokens---before producing a final answer.  The
RL training process (via GRPO or related methods) teaches the model to
\emph{allocate its own test-time compute} by learning when to think longer,
self-verify, backtrack, and explore alternative solution paths.  In this view,
scaling RL training and scaling test-time compute are deeply intertwined: RL
training produces policies that can effectively utilise additional inference-time
tokens, and the quality of this learned compute allocation improves with both
model scale and RL training investment.

%% ---------- Reward Hacking and Overoptimisation ----------
\subsubsection{Reward Hacking and Overoptimisation}

Scaling RL training is not without risks.  Gao et al.\ provided the first
systematic characterisation of \emph{reward model overoptimisation}, the
phenomenon whereby optimising too aggressively against a proxy reward model
degrades true (gold-standard) performance, in accordance with Goodhart's
law~\cite{gao2023scaling}.  Using a synthetic setup in which a large
``gold-standard'' reward model serves as the ground truth, they studied how gold
reward scores evolve as a policy is optimised against a smaller proxy reward
model via both RL and best-of-$n$ sampling.

For best-of-$n$ sampling, the gold reward initially increases as a function of
the KL divergence $d_{\mathrm{KL}}$ between the optimised and initial policies
but eventually peaks and declines.  For RL-based optimisation, a qualitatively
similar pattern emerges, but RL consumes substantially more KL budget for the
same degree of optimisation.  Crucially, the coefficients of these functional
forms \emph{scale smoothly with the number of reward model parameters}: larger
proxy reward models delay the onset of overoptimisation and achieve higher peak
gold rewards, but do not eliminate the phenomenon entirely.  Increasing the
reward model's training dataset size produces a similar protective effect.
The KL penalty coefficient $\beta$ in the RL objective controls how far the
policy drifts from the reference but, notably, does not change the
KL--gold-reward frontier---it merely determines where on the frontier the
optimised policy lands.

These findings have direct implications for GRPO-based training.  Although the
present study uses \emph{rule-based verifiable rewards} (exact-match correctness
on mathematics benchmarks) rather than learned reward models, the
overoptimisation phenomenon can still manifest in subtler forms: the policy may
learn formatting tricks that increase exact-match rates without genuine
reasoning improvement, or it may overfit to the distributional characteristics
of the training prompts.  The KL regularisation term in the GRPO objective
(Section~\ref{sec:grpo}) serves as the primary safeguard against such
degradation, and its coefficient $\beta$ must be tuned with awareness of the
overoptimisation dynamics characterised by Gao et al.

%% ---------- Implications for This Study ----------
\subsubsection{Implications for the Present Study}

The scaling laws surveyed above delineate several orthogonal dimensions along
which performance can be improved: pretraining model size, pretraining data
volume, RL training compute, group size $G$, training data diversity, and
test-time compute.  The present capstone study focuses on a specific and
practically important slice of this space: \emph{the effect of base model size
on GRPO training dynamics and downstream mathematical reasoning performance},
holding the RL algorithm (GRPO), adaptation method (LoRA), training framework
(TRL), reward structure (rule-based verification), and evaluation benchmarks
(GSM8K, MATH) constant.

This design isolates what the scaling literature identifies as one of the most
consequential variables.  Khatri et al.\ showed that scaling model size
substantially raises the asymptotic reward ceiling $A$ in
Equation~\eqref{eq:rl-sigmoid}, while DeepSeek-R1's distillation findings
suggest that there may exist a critical model capacity below which direct RL
training is dominated by knowledge distillation.  By systematically varying
model size while controlling for other factors, this study aims to empirically
locate this transition, characterise how GRPO reward trajectories change with
scale, and provide actionable guidance for practitioners deciding between direct
RL training and distillation under resource constraints.  Additionally, by
monitoring chain-of-thought length and reasoning pattern emergence across
scales, the study connects to the test-time compute scaling literature: if
larger RL-trained models generate longer and more effective reasoning traces,
this would constitute evidence that model scale and learned test-time compute
allocation are complementary rather than substitutable resources in the
mathematical reasoning domain.
